\section{Plan de validación}

\subsection{Estado de la evidencia de mercado}

Al cierre de este documento no se ha aplicado una encuesta ni una entrevista
formal. Por lo tanto, no se presentan porcentajes, tamaños de mercado ni
afirmaciones de disposición de pago. El instrumento está listo y los resultados
se incorporarán cuando se obtengan 10 respuestas de responsables de gimnasios y
10 de miembros adultos.

\begin{table}[H]
\caption{Matriz pendiente de resultados de encuesta}
\begin{tabularx}{\textwidth}{>{\bfseries}p{4cm}p{3cm}X}
\toprule
Indicador & Resultado & Uso esperado \\
\midrule
Método administrativo predominante & \pendiente & Confirmar o rechazar la hipótesis de dispersión. \\
Problemas reportados con mayor frecuencia & \pendiente & Priorizar propuesta de valor. \\
Funciones prioritarias para responsables & \pendiente & Definir alcance del plan básico. \\
Disposición de realizar un piloto & \pendiente & Seleccionar posibles aliados. \\
Rango de pago aceptable & \pendiente & Ajustar precios y equilibrio. \\
Funciones valoradas por miembros & \pendiente & Revisar experiencia y adopción. \\
Preocupaciones de privacidad & \pendiente & Fortalecer incorporación y consentimiento. \\
\bottomrule
\end{tabularx}
\caption*{\textit{Nota}. No existen resultados al 1 de julio de 2026.}
\end{table}

\subsection{Diseño del piloto}

Después de analizar la encuesta se seleccionará un gimnasio que cumpla condiciones
de acceso, interés y responsable designado. Se propone un piloto de cuatro semanas
con datos ficticios durante la capacitación y datos reales únicamente después de
formalizar consentimiento, seguridad y responsabilidades.

\begin{table}[H]
\caption{Indicadores propuestos para un piloto}
\begin{tabularx}{\textwidth}{>{\bfseries}p{4cm}X p{3cm}}
\toprule
Indicador & Método & Meta inicial \\
\midrule
Activación & Responsables que completan configuración. & Definir con aliado. \\
Uso semanal & Sesiones activas por rol. & Establecer línea base. \\
Tareas completadas & Pagos, membresías y asistencias registradas. & Comparar con proceso previo. \\
Incidencias & Errores o dudas por semana. & Tendencia descendente. \\
Tiempo administrativo & Minutos para tareas seleccionadas. & Medir antes y después. \\
Continuidad & Intención de seguir y pagar. & Validar al cierre. \\
\bottomrule
\end{tabularx}
\caption*{\fuentepropia}
\end{table}

\section{Control de requisitos ExpoTÉCNICA}

Los lineamientos definen el modelo de negocio como producto de la etapa
institucional y exigen coherencia entre documento, evidencia y exposición
\parencite{mep2026}. La Tabla~\ref{tab:formularios} registra los controles que el
equipo debe confirmar con el Comité Técnico de Revisión (CTR).

\begin{longtable}{>{\bfseries}p{2.4cm}p{6cm}p{5.6cm}}
\caption{Control de formularios ExpoTEC}\label{tab:formularios}\\
\toprule
Formulario & Uso & Acción para GymHub \\
\midrule
\endfirsthead
\toprule
Formulario & Uso & Acción para GymHub \\
\midrule
\endhead
ExpoTEC-1 & Inscripción del proyecto. & Completar en cada etapa. \\
ExpoTEC-2 & Consentimiento para fotografías o imágenes. & Aplicar si se usan imágenes identificables. \\
ExpoTEC-5 & Experimentos sociales con seres humanos. & Consultar al CTR por encuesta y pruebas. \\
ExpoTEC-6 & Consentimiento informado. & Consultar al CTR antes de recolectar respuestas o pruebas de uso. \\
ExpoTEC-8 & Evaluación de exposición del modelo. & Usar para preparar defensa institucional. \\
ExpoTEC-11 & Evaluación del escrito del modelo. & Ejecutar revisión final del documento. \\
ExpoTEC-9 y 12 & Evaluación regional/nacional. & Reservar si el proyecto avanza de etapa. \\
\bottomrule
\caption*{\textit{Nota}. Síntesis elaborada a partir de los lineamientos ExpoTÉCNICA 2026 \parencite{mep2026}.}
\end{longtable}

\subsection{Privacidad, imagen y propiedad intelectual}

La encuesta no solicitará nombres, correos ni datos de salud. Se limitará a
personas adultas e informará finalidad educativa, voluntariedad y uso agregado.
Las capturas ocultarán identificadores y utilizarán cuentas de demostración. La
base legal y los procedimientos definitivos deberán revisarse antes de operar con
clientes, conforme a la Ley n.º 8968 \parencite{ley8968}.

Las imágenes serán propias, licenciadas o citadas. No se usarán logos de gimnasios,
personas identificables ni material de terceros sin permiso. El Canvas se acredita
como elaboración del equipo.

\subsection{Preparación de la exposición}

La defensa debe poder explicar problema, solución, cliente, Canvas, tecnología,
costos, evidencia y limitaciones en 15 minutos. Se llevará una demostración con
datos ficticios, video offline, capturas, laptop cargada, cargador y extensión
rotulada si la organización lo permite. Ambos estudiantes deberán responder sobre
el negocio y la tecnología.

\subsection{Declaración del estado del proyecto}

GymHub no se presenta como proyecto de continuación, sino como propuesta
desarrollada para ExpoTÉCNICA 2026. Tampoco se describe como producto comercial
terminado. Es un MVP funcional con despliegue demostrativo y trabajo pendiente de
validación, infraestructura y formalización.
